\documentclass[10pt,a4paper]{amsart}
\usepackage[margin=19mm]{geometry}
\usepackage{amsmath,amssymb,mathtools}
\usepackage[scr=rsfs,cal=euler]{mathalfa}
\usepackage{xfrac}
\usepackage[unicode,hidelinks,bookmarks=false]{hyperref}
\newcommand{\ie}{i.e.\ }
\newcommand{\cf}{cf.\ }
\newcommand{\resp}{respectively\ }
\newcommand{\Subsection}{\subsection}
\providecommand{\nobreakdash}{\nobreak\hbox{-}}
\setlength{\emergencystretch}{2em}
\multlinegap0pt
\usepackage{amssymb}
\usepackage{enumerate}
\usepackage{float}
\usepackage{mathtools}
\usepackage[normalem]{ulem}
\usepackage{tikz}
\newtheorem{theorem}{Theorem}
\newtheorem{prop}{Proposition}
\title{A linear upper bound on the zero-sum Ramsey number of forests in $\mathbb{Z}_p$}
\author{Lucas Colucci, Marco D'Emidio}
\date{Source: arXiv:2512.06229v2 (CC BY 4.0)}
\begin{document}
\maketitle

\noindent\emph{Source and licence.} A linear upper bound on the zero-sum Ramsey number of forests in $\mathbb{Z}_p$. Version arXiv:2512.06229v2 (CC BY 4.0), distributed by arXiv under the Creative Commons Attribution 4.0 International licence, http://creativecommons.org/licenses/by/4.0/, as recorded in the arXiv licence metadata for this exact version and shown on the abstract page https://arxiv.org/abs/2512.06229v2. Full licence notice: \texttt{licenses/arxiv-2512.06229-CC-BY-4.0.txt}.\par\smallskip

\noindent\textit{Editorial scope.} Counted unit: the article's prop CauDavProf1 (source0.tex lines 250--252). The article's complete original proof of it is delivered with it (source0.tex lines 253--286, one \texttt{proof} environment of 34 lines). No mathematics is written, repaired or re-derived: the statement and the proof are the article's own lines, in the article's own notation, and no retained line is substituted or re-flowed.  Retained uncounted as the source-local context the proof needs: lines 193--196.  Nothing was omitted from any retained span, so the package carries no omission record.  No retained line contains a citation, so the wrapper carries no bibliography and no bibliography entry.  No retained line contains a figure environment, an image-inclusion command or drawing code, so no image is delivered and the package declares no asset dependency.  Editorial formatting only, in addition: an AMS \texttt{amsart} preamble that restates the article's own declarations and macros used by the retained lines, namely the environments \texttt{theorem}, \texttt{prop} on separate counters, together with the standard packages those lines need.  The retained environments print in this document as Theorem 1, Proposition 1 in order of appearance, whereas the article prints the same items with the numbering of its own full document.  The complete native source of the article as distributed in the arXiv e-print for this exact version is bundled unchanged as \texttt{sources/190474/source0.tex}.
\begin{theorem}[Generalized Cauchy-Davenport] \label{thm:CauDav}
    If $A_1,A_2,\cdots,A_n$ are nonempty subsets of $\mathbb{Z}_p$, then
    $$|A_1+A_2+\cdots+A_n| \geq \min(p,|A_1|+|A_2|+\cdots+|A_n|-n+1).$$
\end{theorem}

\begin{prop}\label{CauDavProf1}
    One of these embeddings of $F$ is zero-sum.
\end{prop}

\begin{proof}
    For each $v_i = u_i$, let $A_i$ be the set of all the possible values of the sum of the edges joining $u_i$ and its $a_i$ selected children when they vary in $X_i\cup Y_i$.

If $X_i = \{p_{i,1},p_{i,2},\cdots,p_{i,a_i}\}$ and $Y_i = \{q_{i,1},q_{i,2},\cdots,q_{i,a_i}\}$ such that the edges $p_{i,j}u_i$ have color $x_i$ for each $j$ and $q_{i,j}u_i$ have color in $S_i$ for each $j$. 

Now, consider the $a_i$ pairs $\mathcal{P}_{i,j} = (p_{i,j},q_{i,j})$ for $j =1,2,\cdots,a_i$. On one hand, by selecting one of $u_i$'s children as one of the elements in $\mathcal{P}_{i,j}$ for each $j$, we obtain $2^{a_j}$ choices of how to embed the children whose possible sums of the $a_i$ edges connecting them to $u_i$ lie on the set $\mathcal{P}_{i,1}+\mathcal{P}_{i,2}+\cdots+\mathcal{P}_{i,a_i}$, meaning:
$$\mathcal{P}_{i,1}+\mathcal{P}_{i,2}+\cdots+\mathcal{P}_{i,a_i}\subseteq A_i$$

On the other hand, since $|\mathcal{P}_{i,j}| = 2$, we can apply the generalized version of the Cauchy-Davenport (Theorem $\ref{thm:CauDav}$):
\begin{align*}
    |A_i|
    \geq |\mathcal{P}_{i,1}+\mathcal{P}_{i,2}+\cdots+\mathcal{P}_{i,a_i}| 
    &\geq \min(p,|\mathcal{P}_{i,1}|+|\mathcal{P}_{i,2}|+\cdots+|\mathcal{P}_{i,a_i}|-(a_i-1))\\
    &\geq \min (p,2+2+\cdots2-(a_i-1))\\
    &\geq \min(p,a_i+1) = a_i+1.
\end{align*}


\medskip

Now, the set of all the possible values of the sum of the edges in the embeddings of $F$ we described earlier is precisely the set $$A \coloneq A_1+A_2+\cdots+A_{m}.$$ 


Therefore, we again apply theorem $\ref{thm:CauDav}$ to bound its size by
\begin{align*}
    |A| 
    &\geq \min(p,|A_1|+\cdots+|A_m|-(m-1))\\
    &\geq\min(p,(a_1+1)+\cdots+(a_m+1)-m+1)\\
    &= \min(p,a_1+\cdots+a_m)\\
    &= \min(p,p) = p.
\end{align*}
    
Proving $A$ spans over all the residues modulo $p$, which implies that $0\in A$ and so there must be a embedding of $T$ that has sum $0$.
\end{proof}

\end{document}
